\scanpage{33}
\mspage{16}
% archive foliation: 32

6. \uline{Relations avec l'octaèdre non orienté} $\Delta_M$.

Soit
\[
(29)\qquad
M^*=\{x\in M\mid x^0\ (\text{image dans }M_0)\ne0\}
=\text{image inverse de }J=M_0^*\subset M_0\ \text{par }M\longrightarrow M_0.
\]
\deleted{\unsure{donc un}}
Donc $M^*$ est un torseur relatif sur $M_0^*=J$, de groupe $(M_0)_J$ (i.e. famille constante de groupes isom. à $M_0$, indexée par $J$). Le groupe $(\mathbb{Z}/4\mathbb{Z})^*=\{\pm1\}$ opère sur $M^*$, de façon compatible avec la projection $M^*\deleted{\unsure{/b}}\longrightarrow J$, on pose
\[
(30)\qquad
\Delta_M=M^*/(\mathbb{Z}/4\mathbb{Z})^*=M^*/\{\pm1\}
=\mathbb{P}(M)(\mathbb{Z}/4\mathbb{Z})
\]
ce sont les points à val. dans $\Lambda=\mathbb{Z}/4\mathbb{Z}$ de la droite projective sur $\Lambda$ définie par $M$, et l'application \deleted{\unsure{unique}}\inserted{can.}
\[
(31)\qquad \Delta_M\longrightarrow J
\]
s'identifie à l'application de restriction, ou réduction mod $2$
\[
\mathbb{P}(M)(\mathbb{Z}/4\mathbb{Z})\longrightarrow\mathbb{P}(M)(\mathbb{F}_2).
\]

C'est une application de degré $2$, qui définit donc sur $\Delta$ la structure de (l'ens. des sommets) d'un octaèdre, dont $J$ est l'ens des diagonales. Les
\deleted{le groupe $GL(M)$ opère sur}
$12$ arêtes de l'octaèdre sont les éléments de $\mathcal P_2(\Delta)$ (qui a $\frac{6\cdot5}{2}=15$ éléments) qui ne sont pas \deleted{des diago} l'une des trois diagonales, i.e. des trois fibres de (31). Les $8$ faces de l'octaèdre sont les parties à trois éléments \inserted{$f$}\deleted{\unsure{$g$}} de $\Delta$ telles que l'application induite par (31) soit une bijection
\[
(32)\qquad f\xrightarrow{\sim}J\ ;
\]
elles s'identifient aux sections de (31), tandis que les arêtes s'identifient aux sections partielles sur des parties de $J$ à deux éléments, et \unsure{bien}
